% !TEX root =./ECE448_Practice_main.tex
%
% L9 practice set (PAM Theory). Every question is a choice we make in the L9
% flowgraphs or a display we read: a symbol rate, an RRC parameter, a Gray
% code, an eye, a loop setting. No error-rate math (ROADMAP, "Sources and
% noise", 2026-10-21): noise is something to watch. Blank copy: \begin{solution}[h] reserves h of work space. SOLUTIONS
% copy (\iskey=1) prints the red worked solutions.

\fullwidth{\textbf{Learning Objective 4.2: } I can describe PAM theory and identify the
pulse-shaping, symbol-rate, and bandwidth trade-offs.
\\[2pt]\textbf{Learning Objective 4.3: } I can apply matched filtering and symbol-timing
recovery to a received PAM signal.
\\[4pt]\normalfont\small Show your work. Write a \textbf{Documentation} line at the
end (\textbf{None} if you did not collaborate).}

\begin{questions}

%%%%%%%%%%%%%%%%%%%%%%%%%%%%%%%%%%%%%%%%%
\question \textbf{Constellations and Gray coding.}
\begin{parts}

	\part Sketch the 4-PAM constellation on the I/Q plane and label its four points
	with a Gray code. Explain why a Gray code is worth using.
		\begin{solution}[1.2in]
		Four points on the real axis at $-3, -1, +1, +3$, labeled 00, 01, 11, 10 from left
		to right. Noise almost always pushes a sample to a neighboring point, and
		neighbors differ in only one bit, so a wrong decision usually costs one bit
		instead of two.
		\end{solution}

	\part 8-PSK places eight points evenly around a circle, three bits each. Label
	them with a Gray code, going around the circle.
		\begin{solution}[1.2in]
		One answer, going counterclockwise from $0^\circ$: 000, 001, 011, 010, 110, 111,
		101, 100. Each point differs from both neighbors in one bit, including the wrap
		from 100 back to 000. Any labeling with that property is correct.
		\end{solution}

	\part On a constellation sink, a 2-PAM signal shows two tight clouds at $\pm 1$.
	The noise is turned up. Describe what changes on the display, and say when the
	receiver starts making wrong decisions.
		\begin{solution}[1.0in]
		Each cloud spreads. While both clouds stay on their own side of the imaginary
		axis, the decision threshold, every decision is right; once the clouds spread far
		enough that some samples cross it, those samples are decided wrong.
		\end{solution}

\end{parts}

\newpage
%%%%%%%%%%%%%%%%%%%%%%%%%%%%%%%%%%%%%%%%%
\question \textbf{Fitting a channel.} A telemetry link must fit a 25 kHz channel and
carry at least 32 kb/s. Use RRC pulse shaping, with bandwidth $B = (1+\alpha)R_s$.
\begin{parts}

	\begin{minipage}[t]{0.58\linewidth}
		\part With 2-PAM and $\alpha = 0.35$, find the largest symbol rate that fits and
		the bit rate it carries. Does it meet the requirement?
		\begin{solution}[1.1in]
		\eq{ R_s = \dfrac{B}{1+\alpha} = \dfrac{25\ \text{kHz}}{1.35} = 18.5\ \text{kbaud} }
		One bit per symbol gives 18.5 kb/s, which is short of 32 kb/s.
		\end{solution}
	\end{minipage}\hfill%
	\begin{minipage}[t]{0.37\linewidth}
		\raggedleft \vspace{12pt}
		$R_s$: \ansbox[1.1in]{18.5 kbaud}

		\vspace{4pt}
		meets it? \ansbox[1.1in]{no}
	\end{minipage}

	\begin{minipage}[t]{0.58\linewidth}
		\part Switch to 4-PAM at the same $\alpha$. Which block setting changes, and what
		is the bit rate now? In words, what would you expect to see on the Symbol Sync
		output at the same noise as before, and why?
		\begin{solution}[1.3in]
		The Chunks to Symbols table changes to four levels, and
		\eq{ R_b = 2 \times 18.5 = 37.0\ \text{kb/s} }
		which meets the requirement. At the same power, four levels sit closer together
		than two, so the same noise spreads each band into its neighbors' thresholds
		sooner: more samples land past a threshold, and more decisions are wrong.
		\end{solution}
	\end{minipage}\hfill%
	\begin{minipage}[t]{0.37\linewidth}
		\raggedleft \vspace{12pt}
		$R_b$: \ansbox[1.1in]{37.0 kb/s}
	\end{minipage}

	\begin{minipage}[t]{0.58\linewidth}
		\part Keep 2-PAM instead and send 32 kbaud. What roll-off would fit 25 kHz?
		Explain why that answer rules out this option.
		\begin{solution}[1.1in]
		\eq{ 1+\alpha = \dfrac{25}{32} = 0.78 \quad\Rightarrow\quad \alpha = -0.22 }
		The roll-off cannot be negative: the narrowest any pulse can make the signal is
		$B = R_s = 32$ kHz at $\alpha = 0$, which is wider than the channel.
		\end{solution}
	\end{minipage}\hfill%
	\begin{minipage}[t]{0.37\linewidth}
		\raggedleft \vspace{12pt}
		\ansbox[1.3in]{impossible}
	\end{minipage}

\end{parts}

\newpage
%%%%%%%%%%%%%%%%%%%%%%%%%%%%%%%%%%%%%%%%%
\question \textbf{Setting up the transmit filter.} A flowgraph runs at
\texttt{samp\_rate} = 48 kHz and sends 2-PAM at 4.8 kbaud with $\alpha = 0.25$.
\begin{parts}

	\begin{minipage}[t]{0.58\linewidth}
		\part Find the samples per symbol and the occupied bandwidth.
		\begin{solution}[1.0in]
		\eq{ \texttt{sps} = \dfrac{48\ \text{kHz}}{4.8\ \text{kbaud}} = 10, \qquad
		B = 1.25 \times 4.8 = 6.0\ \text{kHz} }
		\end{solution}
	\end{minipage}\hfill%
	\begin{minipage}[t]{0.37\linewidth}
		\raggedleft \vspace{12pt}
		\texttt{sps}: \ansbox[1.0in]{10}

		\vspace{4pt}
		$B$: \ansbox[1.0in]{6.0 kHz}
	\end{minipage}

	\part Write the \texttt{firdes.root\_raised\_cosine} call for the transmit taps,
	with a filter 11 symbols long, and name the block that uses them and its
	interpolation.
		\begin{solution}[1.0in]
		\texttt{firdes.root\_raised\_cosine(10, 48000, 4800, 0.25, 110)}, in that order:
		gain \texttt{sps}, sample rate, symbol rate, $\alpha$, and \texttt{11*sps} taps.
		The taps go in an Interpolating FIR Filter with interpolation 10.
		\end{solution}

	\part A teammate wants 4.4 kbaud at the same sample rate. Explain why the
	flowgraph cannot do this as built, and give one way to fix it.
		\begin{solution}[1.0in]
		$48/4.4 = 10.9$ samples per symbol, and the interpolating filter needs an
		integer interpolation. Change the sample rate to an integer multiple of the
		symbol rate, such as $44$ kHz for $\texttt{sps} = 10$, or choose a symbol rate
		that divides 48 kHz, such as 4 kbaud.
		\end{solution}

	\part The teammate shortens the filter to \texttt{2*sps} taps to save computation.
	What happens to the spectrum, and why?
		\begin{solution}[1.0in]
		The filter truncates the pulse to two symbols. Cutting the pulse off abruptly in time
		spreads its spectrum, so the sidelobes outside $\pm B/2$ rise and the signal
		spills into the neighboring channels, which is the problem RRC shaping exists
		to solve.
		\end{solution}

\end{parts}

\newpage
%%%%%%%%%%%%%%%%%%%%%%%%%%%%%%%%%%%%%%%%%
\question \textbf{Reading an eye.} A QT GUI Eye Sink shows a 2-PAM signal at
\texttt{sps} = 8 after the receive filter. The eye is widest open 3 samples after the
start of each symbol period, where the traces cluster at $\pm 0.8$ with a spread of
$\pm 0.2$ around each level.
\begin{parts}

	\begin{minipage}[t]{0.58\linewidth}
		\part Where should the receiver sample, and what is the eye opening there as a
		fraction of the ideal gap from $-0.8$ to $+0.8$?
		\begin{solution}[1.1in]
		Sample 3 samples into each symbol. The innermost traces are at $0.8 - 0.2 = 0.6$
		and $-0.6$, so the opening is
		\eq{ \dfrac{0.6 - (-0.6)}{0.8 - (-0.8)} = \dfrac{1.2}{1.6} = 75\% }
		\end{solution}
	\end{minipage}\hfill%
	\begin{minipage}[t]{0.37\linewidth}
		\raggedleft \vspace{12pt}
		\ansbox[1.1in]{75\%}
	\end{minipage}

	\part For each change below, say how the eye changes: (i) the noise standard deviation doubles;
	(ii) the receiver drops its filter; (iii) the receiver samples 4 samples away from
	the best instant.
		\begin{solution}[1.4in]
		(i) The traces spread vertically, each band gets wider, and the opening between
		them shrinks. (ii) Noise across the whole
		sample bandwidth returns and the single transmit RRC leaves ISI, so the traces
		miss the levels and the eye nearly closes. (iii) Four samples is half a symbol,
		so the receiver samples between symbols. Where the symbol repeats, the sample
		stays near its level, but at a transition it lands near the zero crossing, and
		many of those decisions go wrong.
		\end{solution}

	\part In words, explain why a smaller roll-off makes the receiver more
	sensitive to timing error.
		\begin{solution}[1.0in]
		A smaller $\alpha$ gives a narrower spectrum, which means a longer pulse with
		larger tails. Away from the symbol instant those tails from many neighbors add
		up to more ISI, so the eye's waist is narrower and a given timing error closes
		it more.
		\end{solution}

\end{parts}

\newpage
%%%%%%%%%%%%%%%%%%%%%%%%%%%%%%%%%%%%%%%%%
\question \textbf{The matched filter.} The L9 matched-filter flowgraph runs at
\texttt{samp\_rate} = 32 kHz, \texttt{sps} = 8, and $\alpha = 0.35$.
\begin{parts}

	\part The transmitter uses an RRC filter. What filter should the receiver use,
	and why does that choice also give zero ISI?
		\begin{solution}[1.0in]
		The same RRC filter. The matched filter is the transmit pulse reversed in
		time, and an RRC pulse is symmetric, so it is its own match. The two RRC filters
		in cascade make a raised cosine, which is zero at every other symbol instant.
		\end{solution}

	\part In the L9 matched-filter screenshot, at a noise amplitude of 0.6, the eye
	before the filter is closed and the eye after it is open. The signal occupies about
	$\pm 2.7$ kHz of the 32 kHz the samples represent. Explain in words what the filter
	did to make the difference.
		\begin{solution}[1.0in]
		The noise is spread across the whole 32 kHz, but the signal sits in only a few
		kilohertz of it. The matched filter passes the signal's band and blocks the
		noise everywhere else, so far less noise is left at the sampling instant, the
		traces converge near $\pm 1$, and the eye opens.
		\end{solution}

	\part A classmate puts a full raised-cosine filter at the transmitter and no
	filter at the receiver. The eye at the receiver has zero ISI. Explain what is
	still wrong with this design.
		\begin{solution}[1.0in]
		The receiver passes all the noise across the 32 kHz the samples represent,
		instead of only the signal's band, so the eye is noisier and more decisions are
		wrong at the same transmit power. Zero ISI is necessary but not
		sufficient; the RRC split gets zero ISI and the best noise rejection together.
		\end{solution}

\end{parts}

\newpage
%%%%%%%%%%%%%%%%%%%%%%%%%%%%%%%%%%%%%%%%%
\question \textbf{Timing recovery.} A 2-PAM receiver at 4 kbaud uses Keep 1 in N to
pick one sample per symbol, starting at the center of the eye. Its sample clock is
50 ppm fast relative to the transmitter's.
\begin{parts}

	\begin{minipage}[t]{0.58\linewidth}
		\part How many symbols, and how many seconds, until the sampling point has
		drifted half a symbol, to the zero crossings?
		\begin{solution}[1.1in]
		Each symbol the sampling point moves $50\times10^{-6}$ of a symbol, so
		\eq{ N &= \dfrac{0.5}{50\times10^{-6}} = 10{,}000\ \text{symbols} \\
		t &= \dfrac{10{,}000}{4000} = 2.5\ \text{s} }
		\end{solution}
	\end{minipage}\hfill%
	\begin{minipage}[t]{0.37\linewidth}
		\raggedleft \vspace{12pt}
		symbols: \ansbox[1.1in]{10,000}

		\vspace{4pt}
		time: \ansbox[1.1in]{2.5 s}
	\end{minipage}

	\part You replace Keep 1 in N with a Symbol Sync block. A teammate then puts a
	Keep 1 in N block in front of it to save computation, so Symbol Sync receives one
	sample per symbol, and sets its samples per symbol to 1. Explain why timing
	recovery no longer works, and say where the one-sample-per-symbol step belongs.
		\begin{solution}[1.2in]
		The Gardner detector needs a sample at each symbol instant and one halfway
		between, so it needs at least 2 samples per symbol; with 1 it has no midpoint
		sample and cannot tell early from late.
		Symbol Sync should receive all 8 samples per symbol from the matched filter;
		its own output, at 1 sample per symbol, is the one-sample-per-symbol step.
		\end{solution}

	\part The Symbol Sync output jitters around $\pm 1$ even with little noise. Which
	parameter should change, in which direction, and what is the trade-off?
		\begin{solution}[1.0in]
		Lower the loop bandwidth. A narrower loop averages the timing error over more
		symbols, so the sampling instant jitters less, but the loop locks more slowly
		and follows a drifting clock less closely.
		\end{solution}

\end{parts}

\vspace{1.5em}
\noindent\textbf{Documentation:} \rule[-2pt]{0.6\linewidth}{0.4pt}

\end{questions}
